\documentclass[aspectratio=169,xcolor=table]{beamer}
\usepackage[utf8]{inputenc}
\usepackage[french]{babel}
\usepackage[T1]{fontenc}

\usepackage{booktabs}
\usepackage{subcaption}
\usepackage{graphicx}
\usepackage{colortbl}

\usetheme{Ufg}

% Couleur principale : UFGGreen (autre que UFGBlue) — ou SGITUTeal si votre thème l'accepte
\definecolor{SGITUTeal}{RGB}{0,95,115}
\definecolor{SlideRow}{RGB}{230,245,242}
\definecolor{LightGray}{RGB}{240,240,240}

\setPrimaryColor{UFGGreen}

% Logos désactivés — Université Abdelmalek Essaâdi
% \setLogos{lib/logos/infufgw.png}{lib/logos/infufgw.png}

%------------------------------------- Théorèmes (template)
\newtheorem{conj}{Conjetura}
\newtheorem{defi}{Definição}
\newtheorem{teo}{Teorema}
\newtheorem{lema}{Lema}
\newtheorem{prop}{Proposição}
\newtheorem{cor}{Corolário}
\newtheorem{ex}{Exemplo}
\newtheorem{exer}{Exercício}

\setbeamertemplate{theorems}[numbered]
\setbeamertemplate{caption}[numbered]

%------------------------------------- Document
\begin{document}

\title[SGITU]{SGITU — G4 Coordination des transports}
\subtitle{Support de présentation — Soutenance finale — Pr. BESRI}
\author{EL ASSAL DOUAE, EL MAKOUDI YOUNES, MAGNE TSAFACK LYDIVINE MERVEILLE}
\institute[UAE]{
  Université Abdelmalek Essaâdi\\
  ENSA Tétouan — Génie Informatique 2
}
\date{1\textsuperscript{er} juin 2026}

\frame[noframenumbering]{\titlepage}

%------------------------------------------------ Plan
\setLayout{vertical}
\begin{frame}
  \frametitle{Plan}
  \tableofcontents
\end{frame}

%==============================================================================
\section{Contexte et architecture}
%==============================================================================

\subsection{Contexte SGITU}

\setLayout{vertical}
\begin{frame}{Contexte : où se place G4 ?}
  \begin{itemize}
    \item \textbf{SGITU} : Système de Gestion Intelligente des Transports Urbains
    \item Architecture \textbf{microservices} (G1 à G10)
    \item \textbf{G4} : couche de \textit{coordination opérationnelle}
    \item Périmètre : réseau (lignes, trajets), flotte (missions), événements, intégrations temps réel
  \end{itemize}
  \vspace{0.4cm}
  \begin{center}
    \fbox{\parbox{0.8\textwidth}{\centering
      \textbf{G4} = cerveau opérationnel — port \textbf{8084}\\
      \small G10 (Gateway) $\rightarrow$ JWT $\rightarrow$ API G4
    }}
  \end{center}
\end{frame}

\subsection{Valeur métier}

\begin{frame}{Valeur métier — À quoi sert G4 ?}
  \textbf{Problème :} coordonner l'exploitation quotidienne du réseau.
  \vspace{0.3cm}
  \begin{columns}
    \column{0.48\textwidth}
    \textbf{Gestionnaire réseau (G4\_OPERATOR)}
    \begin{itemize}
      \item Lignes, trajets, arrêts, horaires
    \end{itemize}
    \vspace{0.2cm}
    \textbf{Gestionnaire flotte (DISPATCHER)}
    \begin{itemize}
      \item Missions, affectations
      \item Retards, déviations, pannes
    \end{itemize}
    \column{0.48\textwidth}
    \textbf{Points clés}
    \begin{itemize}
      \item Perturbation \textbf{sans bloquer} la mission
      \item Conflit véhicule $\rightarrow$ HTTP \textbf{409}
      \item Incidents G9 \textbf{séparés} des événements G4
    \end{itemize}
  \end{columns}
  \vspace{0.3cm}
  \begin{block}{Plus-value métier}
    Une perturbation n'arrête pas la mission — on trace un \textbf{événement}.
  \end{block}
\end{frame}

\subsection{Cas d'usage}

\setLayout{vertical}
\begin{frame}{Cas d'usage (exemple soutenance)}
  \textbf{Scénario :} Ligne L12, véhicule 00000000-0000-4000-8000-000000000001 en mission.
  \begin{enumerate}
    \item G7 envoie la position GPS (Kafka \texttt{vehicule-positions})
    \item G4 détecte une \textbf{déviation} $\rightarrow$ événement DEVIATION
    \item G4 notifie G5 (alerte régulation / conducteur)
    \item G9 signale un incident $\rightarrow$ G4 enregistre un \textbf{impact} mission
  \end{enumerate}
  \vspace{0.2cm}
  \small Deux missions \texttt{EN\_COURS} sur le même véhicule $\Rightarrow$ \textbf{409 Conflict} (pas 500).
\end{frame}

\subsection{Architecture technique}

\setLayout{vertical}
\begin{frame}{Architecture technique}
  \begin{itemize}
    \item \textbf{Java 17}, Spring Boot 3.5, Spring Security, Spring Data JPA
    \item \textbf{PostgreSQL} — persistance métier
    \item \textbf{Kafka 3.7} — G7 positions, G9 incidents, G1 lifecycle
    \item \textbf{Docker / Compose} — réseau \texttt{sgitu-network}
    \item \textbf{JWT} aligné G3 / G10
    \item Resilience4j + Prometheus / Grafana
  \end{itemize}
\end{frame}

%==============================================================================
\section{Intégration et API}
%==============================================================================

\subsection{Intégration inter-groupes}

\setLayout{vertical}
\begin{frame}{Intégration inter-groupes}
  \small
  \begin{table}[]
    \centering
    \renewcommand{\arraystretch}{1.3}
    \setlength{\tabcolsep}{8pt}
    {\rowcolors{2}{}{SlideRow}
      \begin{tabular}{p{1.5cm}p{1.5cm}p{5.2cm}}
        \toprule
        \textbf{Groupe} & \textbf{Mode} & \textbf{Contrat} \\
        \midrule
        G3 & HTTP & Validation chauffeurs \texttt{/api/users/drivers/ids} \\
        G5 & HTTP & Notifications + fallback \texttt{DEGRADED} \\
        G7 & Kafka & Topic \texttt{vehicule-positions} \\
        G9 & Kafka & Topic \texttt{incident.transport.topic} \\
        G1 & Kafka & Topic \texttt{missions-lifecycle} \\
        G10 & HTTP & Passerelle JWT \\
        \bottomrule
      \end{tabular}
    }
  \end{table}
\end{frame}

\subsection{Modèle métier}

\begin{frame}{UML — Modèle métier (extrait)}
  \footnotesize
  \begin{itemize}
    \item \textbf{Mission} $1$ — $*$ \textbf{CoordinationEvent} (RETARD, DEVIATION, PANNE…)
    \item \textbf{Mission} $1$ — $0..1$ \textbf{MissionIncidentImpact} $\rightarrow$ référence incident \textbf{G9}
    \item Diagrammes : \texttt{docs/diagrams/*.puml}
  \end{itemize}
  \vspace{0.4cm}
  \begin{block}{Feedback prof}
    Séparation Incident (G9) vs événement opérationnel (G4) — \textbf{intégrée}.
  \end{block}
\end{frame}

\subsection{API REST}

\setLayout{vertical}
\begin{frame}{API REST \& OpenAPI}
  \begin{itemize}
    \item Swagger : \texttt{/swagger-ui.html} — port 8084
    \item Codes HTTP : 200, 201, 400, 401, 403, 404, \textbf{409}, 202
    \item Collection Postman + exemples JSON (\texttt{postman/examples/})
    \item Supervision publique : \texttt{/api/g4/health}, \texttt{/api/g4/logs}
    \item Validation croisée : JWT G3 $\rightarrow$ appel G4
  \end{itemize}
\end{frame}

%==============================================================================
\section{Sécurité et observabilité}
%==============================================================================

\subsection{Sécurité JWT}

\begin{frame}{Sécurité JWT}
  \begin{itemize}
    \item Filtrage \textbf{local} des rôles dans chaque requête (G4 ne fait pas confiance aveugle à G10)
    \item Rôles G3 : \texttt{ROLE\_G4\_OPERATOR}, \texttt{ROLE\_DISPATCHER}, \texttt{ROLE\_G4\_ADMIN}
    \item Validation croisée : token G3 $\rightarrow$ appel G4 (capture Postman)
    \item Secret partagé : \texttt{SGITU\_JWT\_SECRET}
  \end{itemize}
\end{frame}

\subsection{Observabilité}

\setLayout{vertical}
\begin{frame}{Observabilité (pilier livraison)}
  \begin{itemize}
    \item Health enrichi : DB, Kafka, \texttt{pendingNotifications}
    \item Logs structurés : \texttt{group=G4}
    \item \textbf{Prometheus} : \texttt{/actuator/prometheus}
    \item \textbf{Grafana} : dashboard SGITU G4
  \end{itemize}
\end{frame}

\subsection{Résilience}

\begin{frame}{Chaos Monkey — Résilience G5}
  \begin{enumerate}
    \item Prof éteint \textbf{G5}
    \item G4 répond \textbf{202 DEGRADED} (pas de 500)
    \item Notification stockée (\texttt{pending\_notifications})
    \item Renvoi auto quand G5 revient
  \end{enumerate}
  \vspace{0.3cm}
  \begin{block}{Implémentation}
    \small Resilience4j Circuit Breaker + file d'attente locale
  \end{block}
\end{frame}

%==============================================================================
\section{Tests, démo et DevOps}
%==============================================================================

\subsection{Tests et validation}

\setLayout{vertical}
\begin{frame}{Tests \& validation}
  \begin{itemize}
    \item \texttt{./mvnw test} — tests unitaires et intégration (MockMvc)
    \item \texttt{SecurityRolesIntegrationTest} — 403 rôles incorrects
    \item \texttt{MissionVehicleConflictIntegrationTest} — 409 conflit véhicule
    \item Postman : validation croisée + scénarios Chaos Monkey
  \end{itemize}
\end{frame}

\subsection{Démonstration live}

\begin{frame}{Plan de démonstration (soutenance)}
  \textbf{À faire devant le jury (5--7 min) :}
  \begin{enumerate}
    \item \texttt{docker compose --profile monitoring ps}
    \item Postman : login \texttt{gestionnaire.flotte} $\rightarrow$ missions
    \item \texttt{/api/g4/health} et \texttt{/api/g4/logs} (sans token)
    \item \textbf{Bonus :} notification sans G5 sur le réseau $\rightarrow$ statut \texttt{DEGRADED} (Chaos Monkey)
  \end{enumerate}
  \vspace{0.2cm}
  \begin{alertblock}{Conseil}
    Préparer Postman + Docker déjà démarrés avant l'oral.
  \end{alertblock}
\end{frame}

\subsection{DevOps}

\begin{frame}{DevOps — Démarrage}
  \footnotesize
  \texttt{docker compose --profile monitoring up -d --build}
  \vspace{0.3cm}
  \begin{itemize}
    \item Port G4 : \textbf{8084}
    \item Stack : Postgres + Kafka + Prometheus + Grafana
    \item Monorepo : \texttt{SGITU-Microservices/docker-compose.yml}
  \end{itemize}
\end{frame}

\begin{frame}{Conclusion}
  \textbf{G4 livre :}
  \begin{itemize}
    \item Microservice métier \textbf{complet} et \textbf{intégré}
    \item API documentée, sécurisée, conteneurisée
    \item Robustesse démontrable (Chaos Monkey)
  \end{itemize}
\end{frame}

%------------------------------------------------ Slide Merci
\setLayout{blank}
\setBGColor{UFGGreen}
\begin{frame}
  \centering
  \vspace{2cm}
  \textbf{\Huge Merci}
  \\
  \textbf{\Large Questions ?}
  \\
  \vspace{1cm}
  \text{\footnotesize SGITU — G4 Coordination des transports}
  \vspace{2cm}
\end{frame}

\end{document}
